\section{Results}
\label{sec:results}

\subsection{The market beats every model}
\label{sec:results-market}

Every one of the 80 points in Figure~\ref{fig:brier-gap} lies below zero: on
\winB and \winC the market was more accurate than every model in every depth
bucket. The cells share rows across models, so the 80 describe the data rather
than confirm it 80 times; the independent variation is across windows and
buckets.

All 40 \winC cells are clear of zero; on \winB 31 of 40 are, and the nine that
are not form a pattern: the deepest bucket for all eight models, plus
\mGptFiveFive in the thinnest. The market's advantage is smallest in the
deepest bucket, consistent with the depth pattern of
Section~\ref{sec:results-depth}.

The models do beat a naive forecaster. Measured as skill against a forecaster
that always says the base rate, \market scores $+0.199$, $+0.207$ and $+0.232$ on
the three windows, and the roster average scores $+0.121$, $+0.085$ and $+0.072$.
As a counterfactual on the price, pulling the market toward the base rate along the roster's
average fitted line to it scores $+0.158$, $+0.168$ and $+0.194$, still better
than every model; that pull alone reproduces 52, 32 and 24 per cent of the
shortfall. The slope below one, which we call shrinkage, need not mean the
models hedge: a calibrated forecaster that knows less than the market falls back
toward the base rate in the same way, and on \winB and \winC the recalibration
fits of Section~\ref{sec:results-sorting} point away from hedging (on \winA,
where their slopes exceed one, toward it).

Whole-window gaps run from 0.0214 (\mGptFiveFive) to 0.0416 (\mGptOss) on \winB
and from 0.0272 to 0.0431 on \winC, every one clear of zero, as are all four on
\winA (Table~\ref{tab:absolute}). The spread between models is smaller than the
distance from any of them to the market.

Scoring every non-ladder in-play row instead, moved-end-date rows and extreme
prices included, gives Brier edges of $-0.0257$ to $-0.0455$ on \winB and
$-0.0281$ to $-0.0425$ on \winC; ladder rows on their own are worse still on \winB,
$-0.055$ to $-0.067$, and similar on \winC, $-0.024$ to $-0.037$. Every interval
is below zero.

\winA is the one window where the picture is not uniform: the edge is negative in
15 of 20 cells and clear in 8. It is also the smallest window, the only one drawn
without the live-book rule and the only one whose news came from live retrieval
rather than a frozen archive. We report it throughout but rest no claim on it
alone.

Part of the gap is information. News closes about a third of the sorting gap on
\winA and about a sixth on \winB, changing resolution by $-0.002$ to $+0.012$
and raising it for 18 of the 20 model-window pairs. The effect is individually
significant only on \winA, where three of four models have intervals above zero;
on \winB and \winC every per-model interval includes zero, and one model on each
(\mGlmFiveOne, \mDeepSeekVThreeTwo) goes slightly the other way. That \winA is the one window with a significant effect needs a caveat: it is
also the only window whose news was retrieved at run time rather than from a
frozen archive, the setting in which post-cutoff material has been found
leaking through~\citep{ellahib2026}. The gap never
closes: with news in hand, every model still sorts worse than the market. What
these numbers bound is the value of our news (at most three archived articles, on
the minority of rows carrying any), not of information in general; a market
watching lineups, injuries and closing lines has far more.

The gap is several times wider on non-sports rows, so the sports concentration
does not produce it; the models capture a smaller share of what is available on
non-sports than on sports (Appendix~\ref{app:nonsports}). The size of the
deficit depends strongly on forecast horizon, with every model agreeing on the
worst band within a window, but which end of the horizon hurts reverses between
\winB and \winC (Appendix~\ref{app:horizon}).

\mGptFiveFive, one of the two most recently released models, has the highest
resolution in the study and is still short of the market with both intervals
clear of zero; its same-day peer \mDeepSeekVFourPro does not match it
(Appendix~\ref{app:newest}).

\begin{figure*}[!tbp]
\centering
\begin{tikzpicture}
\begin{groupplot}[
  paperaxis,
  group style={group size=2 by 1, horizontal sep=6pt},
  width=0.5\linewidth, height=0.25\linewidth,
  xmin=0.5, xmax=5.5, xtick={1,2,3,4,5},
  xticklabels={1 thin,2,3,4,5 deep},
  xlabel={Depth bucket},
  ymin=-0.11, ymax=0.02,
  ytick={-0.10,-0.08,-0.06,-0.04,-0.02,0},
  scaled y ticks=false,
  yticklabel style={/pgf/number format/fixed, /pgf/number format/precision=2},
  ymajorgrids,
  legend columns=4,
]
\nextgroupplot[title={\winB (\polymarket)}, ylabel={Brier edge (market $-$ model)}, legend to name=legFa,]
\addplot[gray!60, line width=0.4pt, forget plot] coordinates {(0.5,0) (5.5,0)};
\addplot+[mdl/gpt-5.2, forget plot, only marks, mark size=1.5pt, error bars/.cd, y dir=both, y explicit, error bar style={line width=0.5pt}, error mark options={mark size=0pt}] coordinates {(0.685,-0.0268) +- (0,0) -= (0,0.0220) += (0,0.0221) (1.685,-0.0343) +- (0,0) -= (0,0.0223) += (0,0.0215) (2.685,-0.0355) +- (0,0) -= (0,0.0228) += (0,0.0226) (3.685,-0.0209) +- (0,0) -= (0,0.0165) += (0,0.0159) (4.685,-0.0101) +- (0,0) -= (0,0.0150) += (0,0.0134)};
\addplot+[mdl/gpt-oss-120b, forget plot, only marks, mark size=1.5pt, error bars/.cd, y dir=both, y explicit, error bar style={line width=0.5pt}, error mark options={mark size=0pt}] coordinates {(0.775,-0.0447) +- (0,0) -= (0,0.0276) += (0,0.0270) (1.775,-0.0484) +- (0,0) -= (0,0.0220) += (0,0.0215) (2.775,-0.0568) +- (0,0) -= (0,0.0234) += (0,0.0256) (3.775,-0.0410) +- (0,0) -= (0,0.0237) += (0,0.0231) (4.775,-0.0175) +- (0,0) -= (0,0.0230) += (0,0.0195)};
\addplot+[mdl/kimi-k2.5, forget plot, only marks, mark size=1.5pt, error bars/.cd, y dir=both, y explicit, error bar style={line width=0.5pt}, error mark options={mark size=0pt}] coordinates {(0.865,-0.0323) +- (0,0) -= (0,0.0238) += (0,0.0237) (1.865,-0.0408) +- (0,0) -= (0,0.0236) += (0,0.0224) (2.865,-0.0352) +- (0,0) -= (0,0.0197) += (0,0.0198) (3.865,-0.0304) +- (0,0) -= (0,0.0197) += (0,0.0187) (4.865,-0.0083) +- (0,0) -= (0,0.0145) += (0,0.0129)};
\addplot+[mdl/deepseek-v3.2, forget plot, only marks, mark size=1.5pt, error bars/.cd, y dir=both, y explicit, error bar style={line width=0.5pt}, error mark options={mark size=0pt}] coordinates {(0.955,-0.0309) +- (0,0) -= (0,0.0250) += (0,0.0226) (1.955,-0.0403) +- (0,0) -= (0,0.0229) += (0,0.0217) (2.955,-0.0366) +- (0,0) -= (0,0.0214) += (0,0.0214) (3.955,-0.0166) +- (0,0) -= (0,0.0163) += (0,0.0150) (4.955,-0.0119) +- (0,0) -= (0,0.0152) += (0,0.0139)};
\addplot+[mdl/gpt-5.5, forget plot, only marks, mark size=1.5pt, error bars/.cd, y dir=both, y explicit, error bar style={line width=0.5pt}, error mark options={mark size=0pt}] coordinates {(1.045,-0.0156) +- (0,0) -= (0,0.0203) += (0,0.0198) (2.045,-0.0275) +- (0,0) -= (0,0.0205) += (0,0.0189) (3.045,-0.0378) +- (0,0) -= (0,0.0239) += (0,0.0226) (4.045,-0.0185) +- (0,0) -= (0,0.0166) += (0,0.0162) (5.045,-0.0076) +- (0,0) -= (0,0.0147) += (0,0.0130)};
\addplot+[mdl/glm-5.1, forget plot, only marks, mark size=1.5pt, error bars/.cd, y dir=both, y explicit, error bar style={line width=0.5pt}, error mark options={mark size=0pt}] coordinates {(1.135,-0.0321) +- (0,0) -= (0,0.0250) += (0,0.0230) (2.135,-0.0442) +- (0,0) -= (0,0.0252) += (0,0.0231) (3.135,-0.0350) +- (0,0) -= (0,0.0211) += (0,0.0211) (4.135,-0.0214) +- (0,0) -= (0,0.0186) += (0,0.0180) (5.135,-0.0093) +- (0,0) -= (0,0.0165) += (0,0.0137)};
\addplot+[mdl/kimi-k2.6, forget plot, only marks, mark size=1.5pt, error bars/.cd, y dir=both, y explicit, error bar style={line width=0.5pt}, error mark options={mark size=0pt}] coordinates {(1.225,-0.0248) +- (0,0) -= (0,0.0212) += (0,0.0216) (2.225,-0.0419) +- (0,0) -= (0,0.0240) += (0,0.0216) (3.225,-0.0308) +- (0,0) -= (0,0.0221) += (0,0.0210) (4.225,-0.0271) +- (0,0) -= (0,0.0186) += (0,0.0175) (5.225,-0.0073) +- (0,0) -= (0,0.0149) += (0,0.0130)};
\addplot+[mdl/deepseek-v4-pro, forget plot, only marks, mark size=1.5pt, error bars/.cd, y dir=both, y explicit, error bar style={line width=0.5pt}, error mark options={mark size=0pt}] coordinates {(1.315,-0.0359) +- (0,0) -= (0,0.0261) += (0,0.0273) (2.315,-0.0486) +- (0,0) -= (0,0.0250) += (0,0.0251) (3.315,-0.0393) +- (0,0) -= (0,0.0262) += (0,0.0237) (4.315,-0.0195) +- (0,0) -= (0,0.0190) += (0,0.0169) (5.315,-0.0060) +- (0,0) -= (0,0.0175) += (0,0.0148)};
\addlegendimage{mdl/gpt-5.2, only marks, papermark}\addlegendentry{\mGptFiveTwo}
\addlegendimage{mdl/gpt-oss-120b, only marks, papermark}\addlegendentry{\mGptOss}
\addlegendimage{mdl/kimi-k2.5, only marks, papermark}\addlegendentry{\mKimiKTwoFive}
\addlegendimage{mdl/deepseek-v3.2, only marks, papermark}\addlegendentry{\mDeepSeekVThreeTwo}
\addlegendimage{mdl/gpt-5.5, only marks, papermark}\addlegendentry{\mGptFiveFive}
\addlegendimage{mdl/glm-5.1, only marks, papermark}\addlegendentry{\mGlmFiveOne}
\addlegendimage{mdl/kimi-k2.6, only marks, papermark}\addlegendentry{\mKimiKTwoSix}
\addlegendimage{mdl/deepseek-v4-pro, only marks, papermark}\addlegendentry{\mDeepSeekVFourPro}
\nextgroupplot[title={\winC (\kalshi)}, yticklabels={},]
\addplot[gray!60, line width=0.4pt, forget plot] coordinates {(0.5,0) (5.5,0)};
\addplot+[mdl/gpt-5.2, forget plot, only marks, mark size=1.5pt, error bars/.cd, y dir=both, y explicit, error bar style={line width=0.5pt}, error mark options={mark size=0pt}] coordinates {(0.685,-0.0552) +- (0,0) -= (0,0.0221) += (0,0.0220) (1.685,-0.0519) +- (0,0) -= (0,0.0173) += (0,0.0174) (2.685,-0.0302) +- (0,0) -= (0,0.0137) += (0,0.0142) (3.685,-0.0272) +- (0,0) -= (0,0.0154) += (0,0.0145) (4.685,-0.0216) +- (0,0) -= (0,0.0139) += (0,0.0142)};
\addplot+[mdl/gpt-oss-120b, forget plot, only marks, mark size=1.5pt, error bars/.cd, y dir=both, y explicit, error bar style={line width=0.5pt}, error mark options={mark size=0pt}] coordinates {(0.775,-0.0461) +- (0,0) -= (0,0.0210) += (0,0.0204) (1.775,-0.0603) +- (0,0) -= (0,0.0219) += (0,0.0210) (2.775,-0.0309) +- (0,0) -= (0,0.0165) += (0,0.0162) (3.775,-0.0387) +- (0,0) -= (0,0.0176) += (0,0.0174) (4.775,-0.0398) +- (0,0) -= (0,0.0191) += (0,0.0190)};
\addplot+[mdl/kimi-k2.5, forget plot, only marks, mark size=1.5pt, error bars/.cd, y dir=both, y explicit, error bar style={line width=0.5pt}, error mark options={mark size=0pt}] coordinates {(0.865,-0.0488) +- (0,0) -= (0,0.0229) += (0,0.0222) (1.865,-0.0496) +- (0,0) -= (0,0.0191) += (0,0.0188) (2.865,-0.0290) +- (0,0) -= (0,0.0145) += (0,0.0141) (3.865,-0.0285) +- (0,0) -= (0,0.0152) += (0,0.0155) (4.865,-0.0247) +- (0,0) -= (0,0.0139) += (0,0.0142)};
\addplot+[mdl/deepseek-v3.2, forget plot, only marks, mark size=1.5pt, error bars/.cd, y dir=both, y explicit, error bar style={line width=0.5pt}, error mark options={mark size=0pt}] coordinates {(0.955,-0.0516) +- (0,0) -= (0,0.0217) += (0,0.0211) (1.955,-0.0604) +- (0,0) -= (0,0.0213) += (0,0.0194) (2.955,-0.0330) +- (0,0) -= (0,0.0145) += (0,0.0144) (3.955,-0.0274) +- (0,0) -= (0,0.0148) += (0,0.0145) (4.955,-0.0287) +- (0,0) -= (0,0.0147) += (0,0.0148)};
\addplot+[mdl/gpt-5.5, forget plot, only marks, mark size=1.5pt, error bars/.cd, y dir=both, y explicit, error bar style={line width=0.5pt}, error mark options={mark size=0pt}] coordinates {(1.045,-0.0273) +- (0,0) -= (0,0.0180) += (0,0.0181) (2.045,-0.0405) +- (0,0) -= (0,0.0167) += (0,0.0173) (3.045,-0.0241) +- (0,0) -= (0,0.0131) += (0,0.0132) (4.045,-0.0251) +- (0,0) -= (0,0.0138) += (0,0.0140) (5.045,-0.0191) +- (0,0) -= (0,0.0133) += (0,0.0136)};
\addplot+[mdl/glm-5.1, forget plot, only marks, mark size=1.5pt, error bars/.cd, y dir=both, y explicit, error bar style={line width=0.5pt}, error mark options={mark size=0pt}] coordinates {(1.135,-0.0706) +- (0,0) -= (0,0.0252) += (0,0.0251) (2.135,-0.0527) +- (0,0) -= (0,0.0182) += (0,0.0179) (3.135,-0.0291) +- (0,0) -= (0,0.0141) += (0,0.0139) (4.135,-0.0303) +- (0,0) -= (0,0.0154) += (0,0.0146) (5.135,-0.0221) +- (0,0) -= (0,0.0139) += (0,0.0142)};
\addplot+[mdl/kimi-k2.6, forget plot, only marks, mark size=1.5pt, error bars/.cd, y dir=both, y explicit, error bar style={line width=0.5pt}, error mark options={mark size=0pt}] coordinates {(1.225,-0.0477) +- (0,0) -= (0,0.0213) += (0,0.0209) (2.225,-0.0538) +- (0,0) -= (0,0.0188) += (0,0.0188) (3.225,-0.0332) +- (0,0) -= (0,0.0163) += (0,0.0154) (4.225,-0.0361) +- (0,0) -= (0,0.0158) += (0,0.0154) (5.225,-0.0279) +- (0,0) -= (0,0.0137) += (0,0.0139)};
\addplot+[mdl/deepseek-v4-pro, forget plot, only marks, mark size=1.5pt, error bars/.cd, y dir=both, y explicit, error bar style={line width=0.5pt}, error mark options={mark size=0pt}] coordinates {(1.315,-0.0743) +- (0,0) -= (0,0.0281) += (0,0.0265) (2.315,-0.0534) +- (0,0) -= (0,0.0203) += (0,0.0203) (3.315,-0.0301) +- (0,0) -= (0,0.0147) += (0,0.0136) (4.315,-0.0281) +- (0,0) -= (0,0.0171) += (0,0.0164) (5.315,-0.0232) +- (0,0) -= (0,0.0142) += (0,0.0156)};

\end{groupplot}
\end{tikzpicture}

\pgfplotslegendfromname{legFa}

\caption{The market's Brier score minus each model's, per depth bucket; points below zero
mean the market was more accurate. All 80 points lie below zero and 71 of 80 ranges exclude
it; models share rows within a bucket, so the 80 cells are not independent tests. Test split
of the primary set, all rows; whole-window edges with intervals are in Table~\ref{tab:absolute}.}
\label{fig:brier-gap}
\end{figure*}

\subsection{Sorting, not calibration}
\label{sec:results-sorting}

Splitting those Brier scores with Equation~\eqref{eq:murphy} locates the gap.
Figure~\ref{fig:sorting} plots calibration error against resolution: the models
cluster near the market horizontally and far beneath it vertically.

Their reliability runs from 0.0032 to 0.0102 on \winB and from 0.0023 to 0.0106
on \winC, against 0.0017 and 0.0006 for the market. The models are worse here
too, but the quantities are small on both sides, and reliability is a minor part
of the shortfall, not its source.

On the sorting axis the models are not close. Resolution runs from 0.0167 to
0.0306 on \winB and from 0.0194 to 0.0302 on \winC, against 0.050 and 0.056 for
the market, and the market-minus-model
difference is clear of zero for all sixteen model-window pairs ($+0.019$ to
$+0.033$ on \winB, $+0.026$ to $+0.036$ on \winC). By isotonic regression,
resolution accounts for 87 to 99 per cent of each model's Brier gap on \winA, 82
to 97 per cent on \winB and 79 to 95 per cent on \winC. The ten-bin decomposition agrees at 79 to 94 per cent on
\winB and 76 to 94 per cent on \winC; the bin-free figure avoids the within-bin
bias that equal-width bins introduce when the two sides' answers cluster
differently. The pattern holds within both halves of the sample but is weaker on
non-sports rows, where the gap per row is widest: there resolution carries 73 to
90 per cent of the shortfall on \winB and 59 to 89 per cent on \winC, and a
cross-fitted rescaling recovers up to 30 and 42 per cent
(Appendix~\ref{app:nonsports}).

Recalibration is a second view of the same fact, not an independent test, since
reliability is roughly what rescaling can recover (Figure~\ref{fig:recalibration}). The Platt
map recovers 2 to 16 per cent of the gap on \winB and 6 to 18 per cent on \winC;
cross-fitted Platt recovers 1 to 16 per cent and 5 to 15 per cent; a monotone
step function reaches at best 19 and 20 per cent. Four fifths or more is beyond
any monotone rescaling, which cannot reorder forecasts, and it is the ordering
that is wrong. On \winA the fitted maps have slopes above one and recover between
$-9$ and $+2$ per cent. Redrawing the calibration split to keep every event on
one side does not raise recovery (Appendix~\ref{app:splitcheck}).

The Platt slopes (of the recalibration map, not of the fit to the price) on
\winB and \winC are all below one (0.28 to 0.73 on \winB,
0.41 to 0.84 on \winC), meaning the right correction is to pull each forecast
\emph{in}, toward the map's own centre. Yet Figure~\ref{fig:forecast-vs-price} shows the
models answering 0.55 where the crowd says 0.85 and 85 per cent of such rows
resolve Yes. Both are true, of different things: the models are overconfident in
their own noisy ordering (their spread is wider than their ability to
discriminate justifies) and underconfident relative to the market's better
ordering. That reliability comes out small is not evidence of good calibration:
it is a squared quantity set against a resolution gap several times larger, the
worst-calibrated forecasts sit in the sparse tails, and ten bins hide
disagreement inside a bin. Splitting each bin by whether the market priced the
row above or below it raises the measured value three- to fivefold
(Appendix~\ref{app:relsmall}).

Figure~\ref{fig:forecast-vs-price} shows what this means in practice. A
calibrated model's 0.55 is right on average, but the rows behind it can mix
markets the crowd prices at 0.85 with markets it prices at 0.40. The model's 0.55
is honest about the average and silent about the difference, and it is the
difference the Brier score pays for.

\begin{figure*}[!tbp]
\centering
\begin{tikzpicture}
\begin{groupplot}[
  paperaxis,
  group style={group size=2 by 1, horizontal sep=16pt},
  width=0.49\linewidth, height=0.28\linewidth,
  xmin=0, xmax=0.012, ymin=0, ymax=0.06,
  xtick={0,0.004,0.008,0.012}, ytick={0,0.02,0.04,0.06},
  scaled ticks=false,
  xticklabel style={/pgf/number format/fixed, /pgf/number format/precision=3},
  yticklabel style={/pgf/number format/fixed, /pgf/number format/precision=2},
  xlabel={Calibration error (reliability, lower is better)},
  grid=major,
  legend columns=3,
]
\nextgroupplot[title={\winB (\polymarket)}, ylabel={Resolution (sorting, higher is better)}, legend to name=legFb,]
\addplot[mdl/gpt-5.2, forget plot, only marks, papermark] coordinates {(0.0044,0.0268)};
\addplot[mdl/gpt-oss-120b, forget plot, only marks, papermark] coordinates {(0.0102,0.0167)};
\addplot[mdl/kimi-k2.5, forget plot, only marks, papermark] coordinates {(0.0045,0.0229)};
\addplot[mdl/deepseek-v3.2, forget plot, only marks, papermark] coordinates {(0.0032,0.0249)};
\addplot[mdl/gpt-5.5, forget plot, only marks, papermark] coordinates {(0.0036,0.0306)};
\addplot[mdl/glm-5.1, forget plot, only marks, papermark] coordinates {(0.0052,0.0253)};
\addplot[mdl/kimi-k2.6, forget plot, only marks, papermark] coordinates {(0.0042,0.0262)};
\addplot[mdl/deepseek-v4-pro, forget plot, only marks, papermark] coordinates {(0.0071,0.0259)};
\addplot[mdl/market, forget plot, only marks, mark size=3.4pt] coordinates {(0.0017,0.0495)};
\addlegendimage{mdl/gpt-5.2, only marks, papermark}\addlegendentry{\mGptFiveTwo}
\addlegendimage{mdl/gpt-oss-120b, only marks, papermark}\addlegendentry{\mGptOss}
\addlegendimage{mdl/kimi-k2.5, only marks, papermark}\addlegendentry{\mKimiKTwoFive}
\addlegendimage{mdl/deepseek-v3.2, only marks, papermark}\addlegendentry{\mDeepSeekVThreeTwo}
\addlegendimage{mdl/gpt-5.5, only marks, papermark}\addlegendentry{\mGptFiveFive}
\addlegendimage{mdl/glm-5.1, only marks, papermark}\addlegendentry{\mGlmFiveOne}
\addlegendimage{mdl/kimi-k2.6, only marks, papermark}\addlegendentry{\mKimiKTwoSix}
\addlegendimage{mdl/deepseek-v4-pro, only marks, papermark}\addlegendentry{\mDeepSeekVFourPro}
\addlegendimage{mdl/market, only marks, mark size=3.4pt}\addlegendentry{Market price}
\nextgroupplot[title={\winC (\kalshi)}, yticklabels={},]
\addplot[mdl/gpt-5.2, forget plot, only marks, papermark] coordinates {(0.0063,0.0237)};
\addplot[mdl/gpt-oss-120b, forget plot, only marks, papermark] coordinates {(0.0073,0.0194)};
\addplot[mdl/kimi-k2.5, forget plot, only marks, papermark] coordinates {(0.0046,0.0244)};
\addplot[mdl/deepseek-v3.2, forget plot, only marks, papermark] coordinates {(0.0055,0.0206)};
\addplot[mdl/gpt-5.5, forget plot, only marks, papermark] coordinates {(0.0023,0.0302)};
\addplot[mdl/glm-5.1, forget plot, only marks, papermark] coordinates {(0.0066,0.0222)};
\addplot[mdl/kimi-k2.6, forget plot, only marks, papermark] coordinates {(0.0058,0.0216)};
\addplot[mdl/deepseek-v4-pro, forget plot, only marks, papermark] coordinates {(0.0106,0.0244)};
\addplot[mdl/market, forget plot, only marks, mark size=3.4pt] coordinates {(0.0006,0.0557)};

\end{groupplot}
\end{tikzpicture}

\pgfplotslegendfromname{legFb}

\caption{The gap is sorting, not calibration. Each point is one forecaster's calibration error
(across, lower is better) against its resolution (up, higher is better). Ten probability bins;
the bin-free decomposition gives the same picture. Test split of the primary set.}
\label{fig:sorting}
\end{figure*}

\begin{figure*}[!tbp]
\centering
\begin{tikzpicture}
\begin{groupplot}[
  paperaxis,
  group style={group size=2 by 1, horizontal sep=6pt},
  width=0.5\linewidth, height=0.27\linewidth,
  xmin=0.4, xmax=8.6, xtick={1,2,3,4,5,6,7,8},
  xticklabels={{\mGptFiveTwo},{\mGptOss},{\mKimiKTwoFive},{\mDeepSeekVThreeTwo},{\mGptFiveFive},{\mGlmFiveOne},{\mKimiKTwoSix},{\mDeepSeekVFourPro}},
  xticklabel style={rotate=40, anchor=north east, font=\scriptsize},
  ymin=-25, ymax=105, ytick={0,25,50,75,100},
  yticklabel style={/pgf/number format/fixed},
  yticklabel={\pgfmathprintnumber{\tick}\%},
  ymajorgrids, scaled ticks=false,
  legend columns=1,
]
\nextgroupplot[title={\winB (\polymarket)}, ylabel={Share of the Brier gap that rescaling removes}, legend to name=legFf,]
\addplot[papermark, mark=*, color=black, forget plot, only marks] coordinates {(0.76,7.0) (1.76,16.1) (2.76,7.5) (3.76,2.6) (4.76,4.8) (5.76,5.3) (6.76,5.1) (7.76,1.7)};
\addplot[papermark, mark=square, color=black!65, forget plot, only marks] coordinates {(1.00,5.5) (2.00,16.4) (3.00,3.7) (4.00,2.7) (5.00,0.6) (6.00,7.1) (7.00,4.8) (8.00,13.7)};
\addplot[papermark, mark=triangle, color=black!65, forget plot, only marks] coordinates {(1.24,11.6) (2.24,14.9) (3.24,7.5) (4.24,-4.5) (5.24,0.9) (6.24,10.8) (7.24,0.4) (8.24,19.0)};
\addplot[gray!60, line width=0.4pt, forget plot] coordinates {(0.4,0) (8.6,0)};
\addplot[gray!70, dashed, line width=0.6pt, forget plot] coordinates {(0.4,100) (8.6,100)};
\node[font=\scriptsize, anchor=north east, text=black!70] at (axis cs:8.6,98) {gap fully closed};
\addlegendimage{papermark, mark=*, color=black}\addlegendentry{Platt, fit on the calibration split, scored on the test split}
\addlegendimage{papermark, mark=square, color=black!65}\addlegendentry{Platt, cross-fitted over all rows}
\addlegendimage{papermark, mark=triangle, color=black!65}\addlegendentry{Isotonic, cross-fitted over all rows}
\nextgroupplot[title={\winC (\kalshi)}, yticklabels={},]
\addplot[papermark, mark=*, color=black, forget plot, only marks] coordinates {(0.76,12.5) (1.76,13.1) (2.76,9.3) (3.76,8.2) (4.76,5.7) (5.76,12.7) (6.76,11.7) (7.76,17.9)};
\addplot[papermark, mark=square, color=black!65, forget plot, only marks] coordinates {(1.00,10.4) (2.00,11.5) (3.00,7.6) (4.00,5.7) (5.00,4.9) (6.00,11.8) (7.00,9.4) (8.00,15.2)};
\addplot[papermark, mark=triangle, color=black!65, forget plot, only marks] coordinates {(1.24,12.8) (2.24,8.8) (3.24,3.7) (4.24,5.6) (5.24,3.4) (6.24,9.5) (7.24,6.7) (8.24,20.2)};
\addplot[gray!60, line width=0.4pt, forget plot] coordinates {(0.4,0) (8.6,0)};
\addplot[gray!70, dashed, line width=0.6pt, forget plot] coordinates {(0.4,100) (8.6,100)};
\node[font=\scriptsize, anchor=north east, text=black!70] at (axis cs:8.6,98) {gap fully closed};

\end{groupplot}
\end{tikzpicture}

\pgfplotslegendfromname{legFf}

\caption{Rescaling cannot close the gap. Each point is the share of one model's Brier
shortfall removed by an after-the-fact rescaling. Row sets differ by series: the single-split
maps are fitted on the calibration split and scored on the test split, while the two
cross-fitted series score every row. This is the one analysis not reported on the test split
alone; cross-fitting admits no alternative. \winA is not shown: slopes above one, recovery between $-9$ and $+2$ per cent
(single split) and between $-14$ and $+14$ per cent (cross-fitted).}
\label{fig:recalibration}
\end{figure*}

\subsection{Losses sit where a model contradicts the crowd}
\label{sec:results-crowd}

If the models sort poorly, the damage should sit where a model disagrees with the
crowd. Table~\ref{tab:crowd} splits every row by direction. Rows leaning with
the crowd cost little, at a Brier edge of $+0.005$ on \winA, $-0.003$ on \winB
and $-0.011$ on \winC; rows leaning against sit far below, at $-0.038$, $-0.046$
and $-0.052$, and carry 110, 96 and 91 per cent of the net accuracy deficit (over
100 on \winA because there the with-the-crowd rows favour the models). That much
follows from shrinkage alone: the against category is 53, 60 and 66 per cent of
rows, and 66.8, 69.0 and 60.4 per cent of it is the model favouring the market's
outcome with less confidence. The loss sits where the model favours the other
outcome: those rows are 18, 18 and 26 per cent of all rows, score $-0.084$,
$-0.121$ and $-0.104$ against the market, and carry about 80, 78 and 72 per cent
of the net deficit (Appendix~\ref{app:againstdetail}).

In money the models systematically buy longshots~\citep{snowberg2010}:
against-the-crowd trades buy at 0.28 to 0.31 and win 27 to 29 per cent of the
time, returning $-0.172$ per dollar on \winA and $-0.124$ on \winC. On \winB the
money half is not conclusive: neither group's interval excludes zero, though the
accuracy half is clear on all three windows. Fees are about three quarters of
the trading loss on \winB and two fifths on \winC; on \winA, where fees were
close to nothing, the loss is mispricing (Appendix~\ref{app:feesensitivity}).

Figure~\ref{fig:forecast-vs-price} shows where the disagreement lives. In the
highest price band the crowd says 0.846 on \winB and 0.843 on \winC and about 85
per cent of those markets do resolve Yes; the models answer 0.490 to 0.637 and
0.546 to 0.670. Each model's forecast fitted against the price has a slope on the price well below one
(roster mean 0.57 on \winB; individual slopes 0.45 to 0.65), and its fixed point
sits near the window's base rate. The Yes side costs $+0.192$ Brier
on \winA, $+0.082$ on \winB and $+0.075$ on \winC more than the No side, the
pooled interval above zero on each window (Figure~\ref{fig:forecast-vs-price}).

News does not fix this. Scored directly, it is worth $-0.001$ to $+0.017$
in Brier (Appendix~\ref{app:memory}). A toward/away split of the news effect
appears large, but a random-direction placebo reproduces 99 per cent of it on
\winB and 102 per cent on \winC, and overshoots it on \winA: because the price predicts the outcome, any move
landing closer to it scores better whether or not it carried information.

\begin{table*}[!tbp]
\centering
\caption{Leaning with or against the crowd. Rows where the market sits at exactly 0.5, or the
model exactly at the price, have no direction and are in neither group, so shares do not sum
to 100\%; they still count in the net deficit quoted in the text, which is why shares
recomputed from this table alone differ slightly on \winA. Brier edge is the market's score
minus the model's. All models pooled, test split of the primary set; brackets are 95\%
event-resampling ranges. Return averages each trade's own return after fees, so it is not
win share over price paid.}
\label{tab:crowd}
\begingroup\footnotesize\setlength{\tabcolsep}{3pt}%
\begin{tabular}{@{}llrrllrr@{}}
\toprule
Window & Model leans & Rows & Share & Brier edge & Return per \$1 & Price paid & Win share \\
\midrule
A & With & 2{,}236 & 45\% & +0.005 {\scriptsize[$-$0.003, +0.012]} & $-$0.003 {\scriptsize[$-$0.069, +0.053]} & .710 & .707 \\
 & Against & 2{,}627 & 53\% & $-$0.038 {\scriptsize[$-$0.060, $-$0.017]} & $-$0.172 {\scriptsize[$-$0.309, $-$0.017]} & .294 & .267 \\
\addlinespace
B & With & 4{,}536 & 38\% & $-$0.003 {\scriptsize[$-$0.010, +0.004]} & +0.017 {\scriptsize[$-$0.035, +0.069]} & .695 & .716 \\
 & Against & 7{,}051 & 60\% & $-$0.046 {\scriptsize[$-$0.060, $-$0.033]} & $-$0.060 {\scriptsize[$-$0.173, +0.064]} & .282 & .288 \\
\addlinespace
C & With & 5{,}322 & 31\% & $-$0.011 {\scriptsize[$-$0.018, $-$0.004]} & $-$0.043 {\scriptsize[$-$0.094, +0.001]} & .712 & .700 \\
 & Against & 11{,}356 & 66\% & $-$0.052 {\scriptsize[$-$0.062, $-$0.042]} & $-$0.124 {\scriptsize[$-$0.218, $-$0.034]} & .305 & .287 \\
\bottomrule
\end{tabular}
\endgroup

\end{table*}

\subsection{Correlated errors}
\label{sec:results-correlated}

The models do not fail independently. Simulated forecasters that shrink as hard
but share nothing else correlate at only 0.23 to 0.28 in departures from the
price; the real departure correlations are 0.801 on \winA, 0.795 on \winB and
0.775 on \winC (Figure~\ref{fig:correlated} in Appendix~\ref{app:pairdetail}).
Removing each model's own shrinkage and correlating what is left gives 0.750,
0.734 and 0.708, against a simulated value near zero
(Figure~\ref{fig:shrinkage-null}).

That last comparison is weaker than it looks: the simulated forecasters are built
with independent noise, so their residual correlation is near zero by
construction. The comparison worth making is against explanations that do not
require shared judgment. Every model saw the same question and the same news, so
agreement might reflect a common brief. Adding the row's observable features to
the fit (sports flag, topic, horizon, price band, news indicator) barely moves it:
0.711, 0.726 and 0.695, a fall of about five per cent at most. Adding a per-event
intercept, which absorbs everything the markets of one event share whether
measured or not, gives 0.747, 0.662 and 0.694 on the 942, 475 and 511 rows that
belong to events with more than one market. Those rows are not the full test
split, so the comparison that matters is against themselves: price alone gives
0.785, 0.701 and 0.714 on the same rows, so the event intercept removes about
0.04, 0.04 and 0.02. Neither the observable characteristics nor the event itself
accounts for this residual agreement, though neither would capture a shock
specific to one market
at the forecast moment.

A remaining explanation is shared missing information: the price carries facts
no model was given, and forecasters missing the same facts miss the same way. On
news-bearing rows asked with and without the news, adding it lowers the
residual correlation by 0.047, 0.055 and 0.036, clear of zero on \winA and
\winB, but not clearly through what the models share: on \winA each model's own
spread widens while the shared part does not fall, and on \winB and \winC the
shared part falls but not clearly (Appendix~\ref{app:newsagree}). News cannot
carry line-ups or injuries, and the design cannot separate what remains from
shared training or the shared prompt, which we did not vary
(Section~\ref{sec:limitations}). Forecasters given identical inputs
would agree strongly even if each reasoned perfectly, so the result shows that
the models are not independent opinions, not that they share judgment.

A second measure counts how often two models land on the same incorrect side of
even money. On rows the market itself got right, pairs miss together 4.3, 4.0 and
2.7 times the independent rate, against a shrinkage null of 2.0 to 2.2, 1.8 to
1.9 and 1.4 to 1.5; nulls that also share each row's features, or the event,
stay below the real ratio (Appendix~\ref{app:wrongtogether}).

Same-developer pairs sit at both ends of the range, the tightest (\mKimiKTwoFive
with \mKimiKTwoSix, 0.910 on \winB) and the loosest (\mGptFiveFive with \mGptOss,
0.626 on \winC); with only three developers fielding more than one model we
draw no conclusion
(Appendix~\ref{app:pairdetail}).

That the models agree also decides what combining them can do. The arithmetic
mean of all eight ties the best single model on \winB (Brier 0.2077 against
0.2076, both on the rows every model forecast, which is why they differ slightly
from Table~\ref{tab:absolute}) and is worse on \winC (0.2148 against 0.2117);
the best single model is the same whether picked on the test rows or the
calibration split.
Learned weights help but not enough: non-negative least squares with an intercept, cross-fitted by event, reaches
$-0.0111$, $-0.0208$ and $-0.0261$ against the market on the same rows, an
improvement of 0.0018, 0.0010 and 0.0012 on the best single model there, a small fraction of the width of that model's
own interval. Every interval still lies below zero.
The extremizing factor~\citep{baron2014,satopaa2016} is 1.15 on \winA,
0.75 on \winB and 0.85 on \winC: on the two eight-model windows the data asks
for the pooled forecast to be pulled toward even money rather than away from it.

\begin{figure*}[!tbp]
\centering
\begin{tikzpicture}
\begin{groupplot}[
  paperaxis,
  group style={group size=3 by 1, horizontal sep=34pt},
  width=0.36\linewidth, height=0.25\linewidth,
  grid=major,
  scaled ticks=false,
]
\nextgroupplot[title={Every model shrinks},
  xmin=0, xmax=1, ymin=0, ymax=1,
  xtick={0,0.25,0.5,0.75,1}, ytick={0,0.25,0.5,0.75,1},
  xlabel={Market price $q$}, ylabel={Fitted forecast $a + bq$},
  xticklabel style={/pgf/number format/fixed, /pgf/number format/precision=2},
  yticklabel style={/pgf/number format/fixed, /pgf/number format/precision=2},
  legend to name=legFe, legend columns=3,
]
\addplot[gray!70, dashed, line width=0.5pt, forget plot] coordinates {(0,0) (1,1)};
\addplot[mdl/gpt-5.2, forget plot, line width=0.7pt, mark=none] coordinates {(0,0.1322) (1,0.7246)};
\addplot[mdl/gpt-oss-120b, forget plot, line width=0.7pt, mark=none] coordinates {(0,0.1728) (1,0.6275)};
\addplot[mdl/kimi-k2.5, forget plot, line width=0.7pt, mark=none] coordinates {(0,0.1458) (1,0.7017)};
\addplot[mdl/deepseek-v3.2, forget plot, line width=0.7pt, mark=none] coordinates {(0,0.1761) (1,0.7167)};
\addplot[mdl/gpt-5.5, forget plot, line width=0.7pt, mark=none] coordinates {(0,0.1254) (1,0.7453)};
\addplot[mdl/glm-5.1, forget plot, line width=0.7pt, mark=none] coordinates {(0,0.1407) (1,0.7238)};
\addplot[mdl/kimi-k2.6, forget plot, line width=0.7pt, mark=none] coordinates {(0,0.1364) (1,0.7187)};
\addplot[mdl/deepseek-v4-pro, forget plot, line width=0.7pt, mark=none] coordinates {(0,0.1240) (1,0.7741)};
\addlegendimage{papermark, mark=*, color=black}\addlegendentry{Real models}
\addlegendimage{papermark, mark=square, color=black!65}\addlegendentry{Null: same shrinkage, independent noise}
\addlegendimage{papermark, mark=triangle, color=black!65}\addlegendentry{Null: same shrinkage, permuted residuals}
\nextgroupplot[title={Agreement},
  xmin=0.5, xmax=2.5, xtick={1,2}, xticklabels={{$r(p-q)$},{$r$(residual)}},
  ymin=-0.1, ymax=1.0, ytick={0,0.2,0.4,0.6,0.8,1.0},
  ylabel={Correlation between a pair of forecasters},
  yticklabel style={/pgf/number format/fixed, /pgf/number format/precision=1},
]
\addplot[gray!60, line width=0.4pt, forget plot] coordinates {(0.5,0) (2.5,0)};
\addplot+[papermark, mark=*, color=black, forget plot, only marks, error bars/.cd, y dir=both, y explicit, error bar style={line width=0.6pt}, error mark options={mark size=1.2pt}] coordinates {(0.80,0.7950) +- (0,0) -= (0,0.0280) += (0,0.0240) (1.80,0.7340) +- (0,0) -= (0,0.0330) += (0,0.0290)};
\addplot+[papermark, mark=square, color=black!65, forget plot, only marks, error bars/.cd, y dir=both, y explicit, error bar style={line width=0.6pt}, error mark options={mark size=1.2pt}] coordinates {(1.00,0.2770) +- (0,0) -= (0,0.0160) += (0,0.0180) (2.00,0.0000) +- (0,0) -= (0,0.0110) += (0,0.0120)};
\addplot+[papermark, mark=triangle, color=black!65, forget plot, only marks, error bars/.cd, y dir=both, y explicit, error bar style={line width=0.6pt}, error mark options={mark size=1.2pt}] coordinates {(1.20,0.2580) +- (0,0) -= (0,0.0080) += (0,0.0100) (2.20,0.0100) +- (0,0) -= (0,0.0110) += (0,0.0110)};
\nextgroupplot[title={Wrong together},
  xmin=0.5, xmax=2.5, xtick={1,2}, xticklabels={{All rows},{Market right}},
  ymin=0.8, ymax=5.0, ytick={1,2,3,4,5},
  ylabel={Times the independent rate},
]
\addplot[gray!60, line width=0.4pt, forget plot] coordinates {(0.5,1) (2.5,1)};
\addplot+[papermark, mark=*, color=black, forget plot, only marks, error bars/.cd, y dir=both, y explicit, error bar style={line width=0.6pt}, error mark options={mark size=1.2pt}] coordinates {(0.80,2.4650) +- (0,0) -= (0,0.1300) += (0,0.1400) (1.80,4.0360) +- (0,0) -= (0,0.3890) += (0,0.3930)};
\addplot+[papermark, mark=square, color=black!65, forget plot, only marks, error bars/.cd, y dir=both, y explicit, error bar style={line width=0.6pt}, error mark options={mark size=1.2pt}] coordinates {(1.00,1.8400) +- (0,0) -= (0,0.0340) += (0,0.0460) (2.00,1.9290) +- (0,0) -= (0,0.1000) += (0,0.1080)};
\addplot+[papermark, mark=triangle, color=black!65, forget plot, only marks, error bars/.cd, y dir=both, y explicit, error bar style={line width=0.6pt}, error mark options={mark size=1.2pt}] coordinates {(1.20,1.8680) +- (0,0) -= (0,0.0270) += (0,0.0250) (2.20,1.7650) +- (0,0) -= (0,0.0790) += (0,0.0940)};
\end{groupplot}
\end{tikzpicture}

\pgfplotslegendfromname{legFe}

\caption{Shared shrinkage does not account for the agreement. Left: fitted $p = a + bq$ for
each model on \winB; every slope is below one. Middle and right: simulated forecasters
carrying each model's fitted $a$, $b$ and residual spread but independent noise agree on
departures at 0.26 to 0.28 on this window (0.23 to 0.28 across all three) and on residuals at
essentially zero, by construction; the real models reach $r(p-q)=0.795$, residual $r=0.734$ and a
wrong-together ratio of 2.47. Real values carry 95\% event-bootstrap intervals (1,000 resamples, shrinkage
refitted inside each); simulated ranges are 2.5th--97.5th percentiles over 400 draws. Test
split, rows every model forecast.}
\label{fig:shrinkage-null}
\end{figure*}

\subsection{Depth: no effect in returns, some in accuracy}
\label{sec:results-depth}

If thin books hurt the models, losses should be largest where the book is
thinnest. In returns they are not:
the thin-minus-deep return gap is $-0.018$ on \winB and $+0.004$ on \winC, both
intervals spanning zero by a wide margin, no individual model's gap clears zero,
and regressing return on log depth gives no trend clear of zero. In accuracy the
pattern is different. The models' Brier deficit shrinks as the book deepens: its slope on
log depth is $+0.0075$ on \winB and $+0.0072$ on \winC, both intervals above
zero, and the thin-minus-deep gap is $-0.021$ on \winB, its interval including zero, and $-0.027$ on \winC, clear of it. Holding price band or horizon
fixed keeps the slope clear of zero on \winC but not on \winB, and \winA shows no
depth pattern in either measure. Returns are far noisier than Brier scores, so
the two need not agree.

